% !TeX program = pdflatex
% Polityka konfliktów interesów --- Basilisk Systems P.S.A. (polityka, o której mowa w § 10 ust. 8 umowy spółki; część regulaminu Rady Dyrektorów).
% Odesłania pobierane z psa.aux i basilisk_sha.aux (xr-hyper).
\documentclass[11pt,a4paper]{report}
\usepackage[T1]{fontenc}
\usepackage[utf8]{inputenc}
\usepackage[english,polish]{babel}
\usepackage{xr-hyper}
\usepackage{basilisk_i18n}
\usepackage{lmodern}
\usepackage{microtype}
\usepackage{setspace}
\usepackage{parskip}
\usepackage{needspace}
\usepackage{tabularx}
\usepackage{float}
\externaldocument{psa}[psa.pdf]
\externaldocument{basilisk_sha}[basilisk_sha.pdf]

\BusinessPlanCompany{Basilisk Systems}
\BusinessPlanWatermarkText{Basilisk Systems}
\BusinessPlanWatermarkColor{black!6}
\BusinessPlanWatermarkFontSize{38}{46}
\BusinessPlanWatermarkAngle{38}
\BusinessPlanAuthorsText{Basilisk Systems --- Rada Dyrektorów}
\BusinessPlanConfidentialText{Projekt --- dokument wewnętrzny}
\BusinessPlanFooterText{Basilisk Systems --- polityka konfliktów interesów (projekt)}
\BusinessPlanVersionText{Wersja 1.0 --- projekt}

\setstretch{1.1}
\setlist[itemize]{topsep=0.2em,itemsep=0.12em,parsep=0.05em,leftmargin=1.8em}
\setlist[enumerate]{topsep=0.2em,itemsep=0.15em,parsep=0.05em,leftmargin=2.0em}
\setlength{\parindent}{0pt}
\setlength{\parskip}{0.5em}
\setlength{\emergencystretch}{4em}
\renewcommand{\arraystretch}{1.2}
\renewcommand{\thesection}{\arabic{section}}
\captionsetup{font=small,labelfont=bf}

\definecolor{cHR}{HTML}{1F4E79}
\definecolor{cWZ}{HTML}{1F4E79}
\definecolor{cRD}{HTML}{2E7D32}
\definecolor{cDW}{HTML}{8D6E00}
\definecolor{cX}{HTML}{8E2430}
\definecolor{cPR}{HTML}{6A1B9A}
\tikzset{
  >={Stealth[length=2.2mm]},
  every node/.style={font=\small, align=center},
  box/.style={draw, rounded corners=2pt, minimum height=8mm, inner sep=4pt, text width=#1},
  box/.default=3.6cm,
  wz/.style={box=#1, fill=cWZ!12, draw=cWZ},
  wz/.default=3.6cm,
  rd/.style={box=#1, fill=cRD!12, draw=cRD},
  rd/.default=3.6cm,
  dw/.style={box=#1, fill=cDW!14, draw=cDW},
  dw/.default=3.6cm,
  xx/.style={box=#1, fill=cX!10, draw=cX},
  xx/.default=3.6cm,
  q/.style={draw, diamond, aspect=2.4, inner sep=1.5pt, fill=black!4, text width=#1},
  q/.default=3.4cm,
  lbl/.style={font=\scriptsize, fill=white, inner sep=1.5pt},
  line/.style={->, thick},
}
\newtcolorbox{Prawo}{colback=cHR!5, colframe=cHR!70, boxrule=0.6pt, arc=2pt, left=6pt, right=6pt, top=4pt, bottom=4pt, fonttitle=\small\bfseries, title=Podstawa prawna i dokumenty}
\newtcolorbox{Decyzja}{colback=black!3, colframe=black!45, boxrule=0.5pt, arc=2pt, left=6pt, right=6pt, top=4pt, bottom=4pt}
\newtcolorbox{Przyklad}{colback=cPR!5, colframe=cPR!75, boxrule=0.7pt, arc=2pt, left=6pt, right=6pt, top=4pt, bottom=4pt, fonttitle=\small\bfseries, title=Przykład z praktyki Basiliska}

\newcommand{\field}[1]{\textbf{[#1]}}
\newcommand{\US}[1]{US §~#1}
\newcommand{\UA}[1]{UA §~#1}

\hypersetup{
  pdftitle={Basilisk Systems P.S.A. --- polityka konfliktów interesów},
  pdfauthor={Basilisk Systems --- Rada Dyrektorów}
}

\begin{document}

\BusinessPlanTitlePage
  {Basilisk Systems}
  {Wersja 1.0 --- projekt}
  {Polityka konfliktów interesów}
  {Ujawnianie, ocena i rozstrzyganie konfliktów interesów}
  {Część regulaminu Rady Dyrektorów (US § 10 ust. 8) --- dyrektorzy, Założyciele, pracownicy, współpracownicy}
  {1.0}

\tableofcontents
\newpage

% ================================================================
\section{Cel, podstawa i zakres}

Umowa spółki stanowi, że „szczegółowy podział kompetencji, limity wydatków, zasady raportowania i politykę konfliktu interesów określa regulamin Rady Dyrektorów przyjęty jej uchwałą” (\US{\ref{art:rada} ust.~8}). Niniejszy dokument jest tą polityką. Przyjmuje go Rada Dyrektorów uchwałą jako część regulaminu Rady; \textbf{nie wymaga zmiany umowy spółki}.

W małej spółce Założyciele są jednocześnie akcjonariuszami, dyrektorami, autorami technologii, pracownikami i --- czasem --- pożyczkodawcami. Konflikt interesów jest więc stanem normalnym, a nie wyjątkiem. Celem polityki nie jest jego unikanie za wszelką cenę, lecz \textbf{ujawnienie, udokumentowanie i rozstrzygnięcie przez osobę, która konfliktu nie ma}.

\begin{Prawo}
\begin{itemize}
  \item \textbf{Art.~300\textsuperscript{55} §~1 KSH:} „W przypadku sprzeczności interesów spółki z interesami członka organu, jego współmałżonka, krewnych i powinowatych do drugiego stopnia oraz osób, z którymi jest powiązany osobiście, członek organu powinien ujawnić sprzeczność interesów i wstrzymać się od udziału w rozstrzyganiu takich spraw oraz może żądać zaznaczenia tego w protokole.”
  \item \textbf{Art.~300\textsuperscript{55} §~3 KSH} --- zakaz konkurencji dyrektora; umowa spółki zaostrza go (\US{\ref{art:ip} ust.~5}).
  \item \textbf{Art.~300\textsuperscript{79} KSH} --- w umowie i sporze z dyrektorem Spółkę reprezentuje pełnomocnik powołany uchwałą akcjonariuszy, a wobec dyrektora wykonawczego także dyrektor niewykonawczy (\US{\ref{art:reprezentacja} ust.~3}).
  \item \textbf{Art.~300\textsuperscript{60} KSH} --- umowę z firmą audytorską zawiera pełnomocnik powołany uchwałą akcjonariuszy.
  \item \textbf{Umowa spółki:} \US{\ref{art:rada} ust.~6} (głosowanie i kworum przy konflikcie), \US{\ref{art:rada} ust.~7} (transakcje z podmiotami powiązanymi), \US{\ref{art:wz-zastrzezone} ust.~1 lit.~j--k} (pożyczki z dyrektorami, umowy ponad 100\,000 zł rocznie), \US{\ref{art:akcje} ust.~8 lit.~c} (nabycie akcji własnych).
  \item \textbf{Umowa akcjonariuszy:} \UA{\ref{sha:ustanie}} (stwierdzenie ustania zaangażowania bez udziału zainteresowanego), \UA{\ref{sha:konkurencja} ust.~3} (zwolnienie z zakazu konkurencji bez udziału zainteresowanego), \UA{\ref{sha:finansowanie} ust.~2} (pożyczki Założycieli).
\end{itemize}
\end{Prawo}

\textbf{Polityka obejmuje:} dyrektorów, prokurentów, akcjonariuszy pełniących funkcje w Spółce, pracowników, konsultantów i Partnerów (na podstawie klauzul w umowach), dyrektora powołanego przez inwestora i obserwatora, członków komisji antymobbingowej oraz doradców Spółki.

% ================================================================
\section{Definicje}

\begin{description}[leftmargin=1em, style=nextline]
  \item[Konflikt interesów] sytuacja, w której interes osobisty, majątkowy albo zawodowy osoby objętej polityką --- albo osoby jej bliskiej lub podmiotu powiązanego --- może wpływać albo sprawiać wrażenie, że wpływa, na bezstronne wykonywanie obowiązków wobec Spółki. Obejmuje konflikt \emph{rzeczywisty}, \emph{potencjalny} (może powstać w przyszłości) i \emph{pozorny} (osoba trzecia może uznać, że istnieje).
  \item[Osoba bliska] małżonek, krewni i powinowaci do drugiego stopnia, osoba pozostająca we wspólnym pożyciu oraz inne osoby, z którymi dana osoba jest powiązana osobiście (art.~300\textsuperscript{55} §~1 KSH).
  \item[Podmiot powiązany] podmiot, w którym osoba objęta polityką albo osoba jej bliska posiada co najmniej 10\% udziałów lub głosów, jest członkiem organu, prokurentem, wspólnikiem spółki osobowej albo który kontroluje lub na który wywiera znaczący wpływ; do celów podatkowych stosuje się ponadto definicję z art.~11a ust.~1 pkt~4 ustawy o CIT.
  \item[Koordynator] dyrektor niewykonawczy wyznaczony przez Radę do prowadzenia rejestru i oceny konfliktów; gdy konflikt dotyczy Koordynatora --- inny dyrektor niewykonawczy albo Przewodniczący Rady.
\end{description}

% ================================================================
\section{Mapa konfliktów w Basilisku}

\begin{figure}[H]
\centering
\begin{tikzpicture}[node distance=10mm and 8mm]
  \node[rd=3.2cm, minimum height=14mm] (sp) {\textbf{Basilisk P.S.A.}\\interes Spółki};
  \node[wz=3.4cm, above left=10mm and 2mm of sp] (zal) {\textbf{Założyciele-dyrektorzy}\\umowy IP, wynagrodzenia, pożyczki, Cennik};
  \node[wz=3.4cm, above right=10mm and 2mm of sp] (inw) {\textbf{Inwestorzy}\\dyrektor inwestora, spółki portfelowe};
  \node[dw=3.2cm, left=16mm of sp] (rada) {\textbf{Decyzje Rady}\\vesting, odejścia, zgody};
  \node[dw=3.2cm, right=16mm of sp] (sio) {\textbf{Spółki powiązane}\\siostrzane, zależne, ceny};
  \node[xx=3.4cm, below left=10mm and 2mm of sp] (pra) {\textbf{Pracownicy i konsultanci}\\bliscy, praca dodatkowa, własne firmy};
  \node[xx=3.4cm, below right=10mm and 2mm of sp] (kli) {\textbf{Klienci i urzędy}\\prezenty, sektor obronny};
  \draw[<->, thick] (zal) -- (sp); \draw[<->, thick] (inw) -- (sp);
  \draw[<->, thick] (rada) -- (sp); \draw[<->, thick] (sio) -- (sp);
  \draw[<->, thick] (pra) -- (sp); \draw[<->, thick] (kli) -- (sp);
\end{tikzpicture}
\caption{Sześć obszarów konfliktu interesów. \textcolor{cWZ}{\textbf{Niebieski}} --- właściciele; \textcolor{cDW}{\textbf{złoty}} --- decyzje organów i grupa; \textcolor{cX}{\textbf{czerwony}} --- relacje z ludźmi i otoczeniem.}
\end{figure}

\begin{center}
\small
\begin{tabularx}{\linewidth}{@{}>{\raggedright\arraybackslash}p{3.8cm} X >{\raggedright\arraybackslash}p{3.6cm}@{}}
\toprule
\textbf{Sytuacja} & \textbf{Dlaczego to konflikt} & \textbf{Rozstrzyga} \\
\midrule
Umowa przeniesienia praw IP (Umowa nr~1 i nr~2) z Założycielem-dyrektorem & dyrektor po obu stronach umowy & pełnomocnik z uchwały akcjonariuszy (art.~300\textsuperscript{79} KSH) \\
Wynagrodzenie Założyciela za pracę albo według Cennika & ustala je osoba, która je otrzymuje & Rada bez udziału zainteresowanego; dyrektor za funkcję --- WZ \\
Pożyczka Założyciela dla Spółki & oprocentowanie i zabezpieczenie na korzyść pożyczkodawcy & Rada; gdy pożyczkodawca jest dyrektorem --- także WZ 75\% i pełnomocnik \\
Stwierdzenie ustania zaangażowania i rodzaju odejścia (vesting) & od decyzji zależy liczba zwolnionych akcji & Rada bez udziału zainteresowanego (\UA{\ref{sha:ustanie}}) \\
Założyciel inwestuje w firmę z tej samej branży & zakaz konkurencji & Rada bez udziału zainteresowanego (\US{\ref{art:ip} ust.~5}) \\
Dyrektor powołany przez fundusz, który ma spółkę konkurencyjną w portfelu & informacje Spółki mogą trafić do konkurenta & bariera informacyjna zatwierdzona przez Radę \\
Zatrudnienie osoby bliskiej dyrektora albo lidera Ekipy & ocena kandydata nie jest bezstronna & inna osoba prowadzi rekrutację; ostateczna zgoda dyrektora bez konfliktu \\
Pracownik albo konsultant ma udziały w firmie Partnera & wybór i odbiór dostawy & wyłączenie z wyboru i odbioru; Koordynator \\
Prezent albo zaproszenie od klienta z sektora obronnego & ryzyko korupcyjne i utraty zaufania zamawiającego & zasady z rozdziału~\ref{sec:prezenty} \\
Umowa P.S.A. ze spółką siostrzaną Założycieli & cena może przenosić wartość poza Spółkę & Rada; ponad 100\,000 zł rocznie --- WZ 75\% (lit.~k) \\
Członek komisji antymobbingowej zna strony & bezstronność postępowania & zmiana składu (strategia antymobbingowa) \\
\bottomrule
\end{tabularx}
\end{center}

% ================================================================
\section{Procedura}

\begin{figure}[H]
\centering
\begin{tikzpicture}[node distance=6mm and 9mm]
  \node[box=4.2cm, fill=black!8] (uj) {\textbf{1. Ujawnienie}\\niezwłocznie, przed udziałem w sprawie; forma dokumentowa};
  \node[dw=4.2cm, below=of uj] (rej) {\textbf{2. Rejestr}\\Koordynator wpisuje zgłoszenie};
  \node[q=2.9cm, below=of rej] (oc) {\textbf{3. Ocena}\\czy konflikt istnieje i czy da się nim zarządzić?};
  \node[box=2.6cm, left=6mm of oc] (brak) {\textbf{Brak konfliktu}\\adnotacja w rejestrze};
  \node[rd=4.2cm, below=of oc] (sr) {\textbf{4. Środki}\\wyłączenie z decyzji, bariera informacyjna, zgoda właściwego organu, warunki rynkowe};
  \node[xx=2.8cm, right=6mm of oc] (nie) {\textbf{Konflikt nie do zarządzenia}\\rezygnacja z transakcji albo z funkcji};
  \node[rd=4.2cm, below=of sr] (dok) {\textbf{5. Dokumentacja}\\protokół, uzasadnienie warunków rynkowych};
  \node[wz=4.2cm, below=of dok] (prz) {\textbf{6. Przegląd roczny}\\oświadczenia wszystkich dyrektorów i kluczowych osób; raport dla Rady};
  \draw[line] (uj) -- (rej); \draw[line] (rej) -- (oc);
  \draw[line] (oc) -- node[lbl, above]{brak} (brak);
  \draw[line] (oc) -- node[lbl, above]{nie} (nie);
  \draw[line] (oc) -- node[lbl, left]{tak} (sr);
  \draw[line] (sr) -- (dok); \draw[line] (dok) -- (prz);
\end{tikzpicture}
\caption{Procedura postępowania z konfliktem interesów.}
\end{figure}

\begin{enumerate}
  \item \textbf{Ujawnienie.} Osoba objęta polityką ujawnia konflikt niezwłocznie, najpóźniej przed udziałem w dyskusji, głosowaniu, negocjacjach albo odbiorze --- w formie dokumentowej, Koordynatorowi albo na posiedzeniu Rady (wpis do protokołu, art.~300\textsuperscript{55} §~1 KSH). Ujawnia także konflikt pozorny.
  \item \textbf{Rejestr.} Koordynator prowadzi rejestr konfliktów (wzór w Załączniku~B) i przechowuje go z dokumentacją Rady.
  \item \textbf{Ocena.} Koordynator ocenia w ciągu 7 dni, czy konflikt istnieje i jakie środki są właściwe; w sprawach dyrektorów ocenę zatwierdza Rada bez udziału zainteresowanego.
  \item \textbf{Środki:} wyłączenie z dyskusji i głosowania; przekazanie sprawy innej osobie; bariera informacyjna; zgoda organu właściwego według rozdziału~\ref{sec:organy}; potwierdzenie warunków rynkowych (oferty porównawcze, wycena, Cennik); w ostateczności --- rezygnacja z transakcji albo z funkcji.
  \item \textbf{Dokumentacja} --- protokół, uzasadnienie i dowody rynkowości; przechowywanie co najmniej 5 lat.
  \item \textbf{Przegląd roczny} --- w styczniu każdy dyrektor, prokurent, lider Ekipy i Koordynator składa oświadczenie (Załącznik~A); Koordynator przedstawia Radzie raport z rejestru.
\end{enumerate}

% ================================================================
\section{Kto rozstrzyga}
\label{sec:organy}

\subsection*{5.1. Głosowanie w Radzie Dyrektorów}

\begin{figure}[H]
\centering
\begin{tikzpicture}[node distance=7mm and 10mm]
  \node[box=4.2cm, fill=black!8] (d) {\textbf{Dyrektor ma konflikt w sprawie}};
  \node[rd=4.2cm, below=of d] (u) {ujawnia i \textbf{nie głosuje}; wpis do protokołu};
  \node[rd=4.2cm, below=of u] (k) {\textbf{liczy się do kworum}, jeżeli uczestniczy w posiedzeniu albo głosowaniu};
  \node[q=3.6cm, below=of k] (q) {Czy pozostali dyrektorzy mogą podjąć uchwałę?};
  \node[rd=3.4cm, right=10mm of q] (r) {\textbf{Rada} rozstrzyga głosami dyrektorów bez konfliktu};
  \node[wz=4.2cm, below=of q] (w) {\textbf{Walne Zgromadzenie}\\bezwzględna większość głosów oddanych};
  \draw[line] (d) -- (u); \draw[line] (u) -- (k); \draw[line] (k) -- (q);
  \draw[line] (q) -- node[lbl, above]{tak} (r);
  \draw[line] (q) -- node[lbl, left]{nie} (w);
\end{tikzpicture}
\caption{Uchwała Rady przy konflikcie interesów (\US{\ref{art:rada} ust.~6}; art.~300\textsuperscript{55} §~1 i 300\textsuperscript{58} §~2 KSH).}
\end{figure}

\begin{Przyklad}
Rada liczy 3 dyrektorów: A i~B (Założyciele) oraz~C (niewykonawczy). Spółka zawiera z~A umowę o pracę. A nie głosuje, ale uczestniczy w posiedzeniu, więc kworum (co najmniej połowa) jest zachowane. Uchwałę podejmują B i~C. Umowę w imieniu Spółki podpisuje pełnomocnik powołany uchwałą akcjonariuszy albo --- jeżeli A jest dyrektorem wykonawczym --- C na podstawie uchwały dyrektorów niewykonawczych. Gdyby konflikt miał także B (np. umowa z firmą, w której A i~B mają udziały), uchwałę podejmuje Walne Zgromadzenie.
\end{Przyklad}

\subsection*{5.2. Transakcje z osobami i podmiotami powiązanymi}

\begin{figure}[H]
\centering
\begin{tikzpicture}[node distance=7mm and 10mm]
  \node[q=3.6cm] (q1) {Drugą stroną jest dyrektor?};
  \node[wz=4.6cm, right=of q1] (a1) {\textbf{Pełnomocnik} z uchwały akcjonariuszy; przy dyrektorze wykonawczym także dyrektor niewykonawczy (\US{\ref{art:reprezentacja} ust.~3})};
  \node[q=3.6cm, below=of q1] (q2) {Kredyt, pożyczka, poręczenie z dyrektorem?};
  \node[wz=4.6cm, right=of q2] (a2) {\textbf{WZ 75\%} (\US{\ref{art:wz-zastrzezone} ust.~1 lit.~j}) + uchwała Rady + pełnomocnik};
  \node[q=3.6cm, below=of q2] (q3) {Akcjonariusz $\geq$5\%, dyrektor, osoba bliska albo podmiot powiązany?};
  \node[rd=4.6cm, right=of q3] (a3) {\textbf{Uchwała Rady} bez dyrektora z konfliktem (\US{\ref{art:rada} ust.~7})};
  \node[q=3.6cm, below=of q3] (q4) {Ponad 100\,000 zł w roku obrotowym?};
  \node[wz=4.6cm, right=of q4] (a4) {\textbf{+ WZ 75\%} (\US{\ref{art:wz-zastrzezone} ust.~1 lit.~k})};
  \node[dw=4.0cm, below=of q4] (a5) {zwykłe progi kwotowe (budżet, 250\,000 / 500\,000 zł)};
  \draw[line] (q1) -- node[lbl, above]{tak} (a1);
  \draw[line] (q2) -- node[lbl, above]{tak} (a2);
  \draw[line] (q3) -- node[lbl, above]{tak} (a3);
  \draw[line] (q4) -- node[lbl, above]{tak} (a4);
  \draw[line] (q1) -- node[lbl, left]{dalej} (q2);
  \draw[line] (q2) -- node[lbl, left]{dalej} (q3);
  \draw[line] (q3) -- node[lbl, left]{tak} (q4);
  \draw[line] (q4) -- node[lbl, left]{nie} (a5);
  \draw[line] (q3.west) -- ++(-0.5,0) |- node[lbl, pos=0.25, left]{nie} (a5.west);
\end{tikzpicture}
\caption{Kto zatwierdza i kto podpisuje transakcję z osobą albo podmiotem powiązanym. Pytania stosuje się łącznie --- jedna umowa może wymagać kilku zgód.}
\end{figure}

Wyjątki z \US{\ref{art:wz-zastrzezone} ust.~1 lit.~k}: umowy o pracę i o świadczenie usług Założycieli w zakresie zaangażowania oraz umowy przeniesienia praw (Umowy nr~1 i nr~2) i umowy wykonawcze do obowiązku zbycia --- nie wymagają zgody Walnego Zgromadzenia, ale nadal wymagają uchwały Rady bez udziału zainteresowanego i pełnomocnika po stronie Spółki.

\begin{Prawo}
\textbf{Skutki naruszenia.} Umowa z dyrektorem zawarta bez pełnomocnika jest zagrożona nieważnością (art.~300\textsuperscript{79} KSH) --- w przypadku umów IP oznacza to ryzyko braku praw Spółki do technologii. Czynność dokonana bez zgody wymaganej tylko umową spółki jest ważna wobec kontrahenta, ale dyrektor odpowiada wobec Spółki (\US{\ref{art:reprezentacja} ust.~4}). \textbf{Podatki:} transakcje z podmiotami powiązanymi muszą odpowiadać cenom rynkowym (ceny transferowe, art.~11a--11c ustawy o CIT); przy estońskim CIT nadwyżka ponad wartość rynkową i odsetki od pożyczek akcjonariuszy są ukrytym zyskiem opodatkowanym od razu (art.~28m ust.~3).
\end{Prawo}

\subsection*{5.3. Pozostałe osoby}

\begin{center}
\small
\begin{tabularx}{\linewidth}{@{}>{\raggedright\arraybackslash}p{4.2cm} X@{}}
\toprule
\textbf{Kto ma konflikt} & \textbf{Kto rozstrzyga} \\
\midrule
Dyrektor wykonawczy & dyrektorzy niewykonawczy (\US{\ref{art:reprezentacja} ust.~3}) albo Rada bez jego udziału \\
Wszyscy dyrektorzy albo brak kworum bez konfliktu & Walne Zgromadzenie (\US{\ref{art:rada} ust.~6}) \\
Pracownik, lider Ekipy, Opiekun Systemów & przełożony bez konfliktu; dyrektor wykonawczy \\
Konsultant albo Partner & osoba odpowiedzialna za zlecenie; Koordynator \\
Koordynator & inny dyrektor niewykonawczy albo Przewodniczący Rady \\
Akcjonariusz w głosowaniu WZ & akcjonariusz głosuje, chyba że ustawa wyłącza go od głosowania; umowa spółki nie wprowadza dodatkowych wyłączeń (\US{\ref{art:wz} ust.~7}) \\
\bottomrule
\end{tabularx}
\end{center}

% ================================================================
\section{Zakaz konkurencji i działalność dodatkowa}

\begin{itemize}
  \item \textbf{Założyciele i dyrektorzy} nie prowadzą działalności konkurencyjnej ani nie uczestniczą w podmiocie konkurencyjnym bez zgody Rady podjętej bez udziału zainteresowanego; wyjątek --- inwestycje portfelowe do 5\% bez dostępu do informacji poufnych i bez funkcji w organach (\US{\ref{art:ip} ust.~5}). Umowa spółki jest surowsza niż art.~300\textsuperscript{55} §~3 KSH (próg 10\%).
  \item \textbf{Pracownicy} mogą pracować także dla innych podmiotów (art.~26\textsuperscript{1} KP), ale ujawniają współpracę, która może prowadzić do konfliktu interesów albo naruszenia poufności; zakaz konkurencji --- tylko w odrębnej umowie (wzór~C).
  \item \textbf{Konsultanci} --- zakaz konkurencji w umowie (wzór~B) i oświadczenie o konfliktach interesów (Załącznik~B3).
  \item \textbf{Spółka siostrzana Założycieli} --- zgoda Rady, ceny rynkowe, jasny podział praw do technologii (strategia międzynarodowa, rozdział~5).
\end{itemize}

% ================================================================
\section{Bariera informacyjna}

\begin{figure}[H]
\centering
\begin{tikzpicture}[node distance=8mm and 12mm]
  \node[rd=3.4cm] (s) {\textbf{Basilisk P.S.A.}\\Rada Dyrektorów};
  \node[wz=3.4cm, right=24mm of s] (d) {\textbf{Dyrektor inwestora}\\albo obserwator};
  \node[wz=3.4cm, right=12mm of d] (f) {\textbf{Fundusz}};
  \node[xx=3.4cm, below=12mm of f] (k) {\textbf{Spółka portfelowa}\\konkurencyjna};
  \draw[line] (s) -- node[lbl, above]{informacje} (d);
  \draw[line] (d) -- node[lbl, above]{raporty} (f);
  \draw[line, cX, dashed] (f) -- node[lbl, right]{zakaz} (k);
  \node[box=6.6cm, fill=black!3, below=12mm of s.south east, anchor=north west, xshift=-2.2cm] (nie) {\textbf{Nie przekazuje się:} informacji niejawnych, Informacji Kontrolowanych, informacji wymagających zezwolenia eksportowego, szczegółów technicznych i handlowych istotnych dla konkurenta};
\end{tikzpicture}
\caption{Bariera informacyjna dla dyrektora powołanego przez inwestora (\US{\ref{art:ip} ust.~5}; wzór uprawnień indywidualnych, wariant A i~B).}
\end{figure}

Barierę informacyjną zatwierdza Rada Dyrektorów. Określa ona: osoby po stronie inwestora uprawnione do informacji, kategorie informacji, które nie opuszczają Spółki, sposób przekazywania raportów, wyłączenie dyrektora inwestora z punktów obrad dotyczących konkurenta i potwierdzenie przestrzegania bariery raz w roku. Te same zasady stosuje się do spółki siostrzanej Założycieli i do doradców obsługujących konkurentów.

% ================================================================
\section{Prezenty, zaproszenia i korzyści}
\label{sec:prezenty}

\begin{center}
\small
\begin{tabularx}{\linewidth}{@{}>{\raggedright\arraybackslash}p{4.4cm} X@{}}
\toprule
\textbf{Sytuacja} & \textbf{Zasada} \\
\midrule
Funkcjonariusz publiczny, przedstawiciel zamawiającego z sektora obronnego, osoba prowadząca postępowanie o udzielenie zamówienia & \textbf{żadnych prezentów ani korzyści}, z wyjątkiem drobnych materiałów firmowych o znikomej wartości \\
Prezent od kontrahenta albo dla kontrahenta & do \field{200} zł, okazjonalnie, jawnie; powyżej --- zgoda Koordynatora i wpis do rejestru \\
Zaproszenie na wydarzenie branżowe & dopuszczalne, jeżeli ma charakter merytoryczny; koszty podróży i noclegu pokrywa Spółka \\
Korzyść w trakcie negocjacji, przetargu albo odbioru & niedopuszczalna niezależnie od wartości \\
Gotówka, karty podarunkowe, rabaty prywatne & niedopuszczalne \\
\bottomrule
\end{tabularx}
\end{center}

Wręczenie albo przyjęcie korzyści w związku z pełnieniem funkcji publicznej może stanowić przestępstwo przekupstwa albo sprzedajności (art.~228--229 Kodeksu karnego), a w obrocie gospodarczym --- korupcję menedżerską (art.~296a Kodeksu karnego). Konsultanci są związani zasadami antykorupcyjnymi na podstawie umowy (wzór~B).

% ================================================================
\section{Skutki naruszenia polityki}

\begin{itemize}
  \item \textbf{Dyrektor:} odpowiedzialność wobec Spółki za szkodę; podstawa do odwołania przez Walne Zgromadzenie; w przypadku Założyciela --- ocena, czy naruszenie stanowi istotną przyczynę odejścia w rozumieniu \US{\ref{art:obowiazek-zbycia} ust.~5} (lista zamknięta; konflikt interesów sam w sobie się na niej nie znajduje, chyba że jest jednocześnie naruszeniem poufności, IP albo zakazu konkurencji).
  \item \textbf{Pracownik:} kary porządkowe (art.~108 KP), wypowiedzenie albo rozwiązanie umowy (art.~52 KP), odpowiedzialność za szkodę według Kodeksu pracy.
  \item \textbf{Konsultant, Partner:} rozwiązanie umowy z ważnych powodów, kary umowne przewidziane w umowie.
  \item \textbf{Czynność:} ryzyko nieważności (umowa z dyrektorem bez pełnomocnika), konieczność korekty podatkowej (ceny transferowe, ukryte zyski).
\end{itemize}

% ================================================================
\section{Wdrożenie --- czy trzeba zmieniać umowę spółki?}

\textbf{Nie.} Umowa spółki odsyła politykę konfliktu interesów do regulaminu Rady (\US{\ref{art:rada} ust.~8}). Do wdrożenia potrzebne są:
\begin{enumerate}
  \item uchwała Rady przyjmująca politykę jako część regulaminu Rady i wyznaczająca Koordynatora;
  \item oświadczenia początkowe wszystkich dyrektorów, prokurentów i liderów (Załącznik~A);
  \item uchwała akcjonariuszy powołująca pełnomocnika do umów z dyrektorami --- przy zawiązaniu jest nią uchwała nr~4 z pakietu aktu zawiązania;
  \item wpis do rejestru istniejących konfliktów: umowy IP Założycieli, umowy o pracę Założycieli, ewentualne pożyczki i spółki siostrzane;
  \item klauzula o konfliktach interesów w umowach z konsultantami i Partnerami (istnieje we wzorze~B i Załączniku~B3).
\end{enumerate}

% ================================================================
\section*{Załącznik A --- Wzór oświadczenia o konfliktach interesów}
\addcontentsline{toc}{section}{Załącznik A --- Wzór oświadczenia o konfliktach interesów}

Ja, \field{imię, nazwisko, funkcja}, oświadczam, że według stanu na dzień \field{data}:
\begin{enumerate}
  \item posiadam co najmniej 10\% udziałów lub głosów albo pełnię funkcję w organach następujących podmiotów: \field{brak / wykaz};
  \item osoby mi bliskie są zatrudnione w Spółce albo współpracują z nią: \field{brak / wykaz};
  \item ja albo osoby mi bliskie zawarliśmy ze Spółką następujące umowy: \field{brak / wykaz};
  \item prowadzę dodatkową działalność zawodową: \field{brak / opis};
  \item w ciągu ostatnich 12 miesięcy przyjąłem/przyjęłam albo wręczyłem/wręczyłam prezenty o wartości powyżej \field{200} zł: \field{brak / wykaz};
  \item nie są mi znane inne okoliczności, które mogą stanowić konflikt interesów, z wyjątkiem: \field{brak / opis}.
\end{enumerate}
Zobowiązuję się niezwłocznie ujawniać zmiany i nowe konflikty interesów.

\field{miejscowość, data} \hfill \field{podpis}

\section*{Załącznik B --- Wzór rejestru konfliktów interesów}
\addcontentsline{toc}{section}{Załącznik B --- Wzór rejestru konfliktów interesów}

\begin{center}
\small
\begin{tabularx}{\linewidth}{@{}p{0.6cm} >{\raggedright\arraybackslash}p{2.2cm} >{\raggedright\arraybackslash}X >{\raggedright\arraybackslash}p{2.6cm} >{\raggedright\arraybackslash}p{2.6cm} p{1.6cm}@{}}
\toprule
\textbf{Nr} & \textbf{Osoba} & \textbf{Opis konfliktu i sprawy} & \textbf{Środki} & \textbf{Decyzja (organ, data)} & \textbf{Status} \\
\midrule
1 & \field{\ldots} & \field{np. umowa przeniesienia praw IP z dyrektorem} & \field{pełnomocnik; wyłączenie z głosowania} & \field{uchwała nr \ldots} & \field{otwarty / zamknięty} \\
\bottomrule
\end{tabularx}
\end{center}

\vspace{1em}
{\small Polityka nie zmienia umowy spółki ani umowy akcjonariuszy; w razie rozbieżności rozstrzygają te umowy. Kwotę progu prezentów i definicję podmiotu powiązanego Rada może dostosować uchwałą.\par}

\end{document}
